\documentclass[11pt,a4paper]{article}

\usepackage[margin=2.4cm]{geometry}
\usepackage{amsmath,amssymb}
\usepackage{graphicx}
\usepackage{booktabs}
\usepackage{tabularx}
\usepackage{longtable}
\usepackage{multirow}
\usepackage{xcolor}
\usepackage{hyperref}
\usepackage{enumitem}
\usepackage{caption}
\usepackage{subcaption}
\usepackage{float}
\usepackage{listings}
\usepackage{fancyhdr}
\usepackage{titlesec}
\usepackage{microtype}
\usepackage{array}

\graphicspath{{./reports/figures/}{./}}
\hypersetup{
  colorlinks=true,
  linkcolor=blue!60!black,
  urlcolor=blue!70!black,
  pdftitle={QF5214 Aviation Quant Pipeline}
}

\pagestyle{fancy}
\fancyhf{}
\lhead{\footnotesize QF5214 Aviation Quant Pipeline}
\rhead{\footnotesize Group Project Report}
\cfoot{\thepage}
\renewcommand{\headrulewidth}{0.4pt}

\titleformat{\section}{\large\bfseries}{\thesection}{1em}{}[\titlerule]
\titleformat{\subsection}{\normalsize\bfseries}{\thesubsection}{1em}{}
\titleformat{\subsubsection}{\normalsize\itshape}{\thesubsubsection}{1em}{}

\lstset{
  basicstyle=\ttfamily\footnotesize,
  frame=single,
  breaklines=true,
  backgroundcolor=\color{gray!5},
  captionpos=b,
  numbers=left,
  numberstyle=\tiny\color{gray},
  xleftmargin=1.5em
}

\newcommand{\RPK}{\mathrm{RPK}}
\newcommand{\ASK}{\mathrm{ASK}}
\newcommand{\LF}{\mathrm{LF}}
\newcommand{\IC}{\mathrm{IC}}
\newcommand{\WMAPE}{\mathrm{WMAPE}}

\begin{document}

\begin{titlepage}
\centering
\vspace*{1.4cm}
{\LARGE\bfseries
QF5214 Aviation Quant Pipeline:\\[0.4em]
Data Engineering and Data Science for\\[0.4em]
Airline Earnings-Oriented Research\par}
\vspace{1.2cm}
{\large Group Project Report\par}
\vspace{0.8cm}
{\normalsize April 2026\par}
\vspace{1.4cm}

\begin{abstract}
\noindent
This report presents QF5214, an end-to-end aviation quantitative research
pipeline that starts from raw BTS T-100 segment files and ends with
interpretable backtest results, figures, and automated reports. The project is
designed around two priorities. The first is data engineering: organizing large
raw airline operating records into a reproducible SQLite-centered workflow with
stable intermediate tables, cache-aware execution, and clear output artifacts.
The second is data science: transforming route-level airline operations into
model-ready features, predicting load factor and passenger traffic out of
sample, aggregating the resulting signals into quarterly factors, and testing
whether those factors are informative in an event-driven equity framework.

The implemented pipeline consists of data cleaning, feature engineering,
XGBoost-based load-factor modeling, out-of-sample stochastic prediction,
quarterly factor construction, and a six-airline event-driven backtest using a
multiplicative-weights-update combination rule. In its current checked state,
the prediction layer achieves $\WMAPE \approx 11.34\%$ over 2,656,668
out-of-sample observations. The final backtest covers 31 quarters and produces
positive gross and net cumulative returns of approximately $91.47\%$ and
$68.15\%$, respectively, with annualized Sharpe ratios of 1.287 gross and
1.027 net.

The report argues that the strongest contribution of the project lies in the
integration of engineering and modeling rather than in any single metric. The
system demonstrates how raw airline operational data can be structured,
transformed, modeled, and evaluated inside one coherent workflow. This makes
the project suitable not only as a quantitative analysis exercise but also as a
coursework example of how data engineering and data science can be combined in
a practical research setting.
\end{abstract}

\vfill
\begin{tabular}{ll}
\textbf{Primary data source:} & BTS T-100 Segment data \\
\textbf{Core storage layer:} & SQLite database \texttt{data/quant\_flights.db} \\
\textbf{Main model:} & XGBoost load-factor regressor \\
\textbf{Main evaluation:} & Quarterly event-driven airline equity backtest \\
\end{tabular}
\end{titlepage}

\tableofcontents
\newpage

\section{Introduction and Problem Statement}

\subsection{Project Motivation}

Airline companies are unusual in that a large part of their underlying business
activity is visible through operational data before the market sees the final
quarterly report in structured financial form. Passenger counts, seat supply,
route activity, and traffic volume all evolve during the quarter, while the
earnings release is published only after the operating period is complete. This
creates a natural research question: can operational airline data be organized
early enough, and modeled carefully enough, to provide useful information about
later earnings-related market outcomes?

The motivation for this project comes directly from that timing gap. Rather than
starting from analyst estimates or purely market-based factors, the project
starts from the operating layer. This is important because airline economics
are strongly tied to tangible quantities such as passengers carried, distance
flown, and available seat supply. A pipeline that can structure these data
cleanly may therefore support both business interpretation and quantitative
evaluation.

\subsection{Problem Statement}

The project addresses a linked set of engineering and modeling problems.

From the \textbf{data engineering} perspective, the problem is how to ingest
large raw BTS route-level files, normalize them into consistent tables, and
preserve intermediate outputs in a form that can be reused across pipeline
stages. The challenge is not only volume but also workflow structure: without a
clear storage design, the feature, model, and backtest stages would become hard
to reproduce or explain.

From the \textbf{data science} perspective, the problem is how to translate
route-level operational data into signals that are useful at the quarterly
carrier level. This requires feature engineering, predictive modeling,
aggregation, and event-based evaluation. The modeling goal is not only to fit a
route-level forecast but also to test whether the resulting signals can support
an interpretable equity-event research framework.

\subsection{Project Objective}

The objective of QF5214 is to build one complete workflow that:

\begin{enumerate}[leftmargin=*, label=(\roman*)]
  \item ingests and stores raw airline operating data efficiently;
  \item constructs route-level features from those operating records;
  \item trains an out-of-sample predictive model for airline load factor and
  related passenger quantities;
  \item converts model outputs into quarterly factors at the airline level; and
  \item evaluates those factors in a tradable event-driven backtest around
  earnings releases.
\end{enumerate}

This objective is deliberately broader than a single-model benchmark. The
report treats the project as a pipeline, not as a one-off model experiment.
That framing is essential because the strongest value of the project comes from
the interaction between its storage design, feature logic, predictive model,
and final evaluation layer.

\subsection{Report Orientation}

This report is written with a clear emphasis on the two most important
technical layers of the project:

\begin{itemize}[leftmargin=*]
  \item \textbf{Data engineering}, including ingestion, storage design,
  intermediate tables, caching, and reproducibility; and
  \item \textbf{Data science}, including feature construction, predictive
  modeling, factor construction, and event-driven backtest interpretation.
\end{itemize}

The project does contain a lightweight automation layer for current-run report
generation and anomaly summarization, but that is intentionally treated as a
supporting component rather than the center of the report.

\section{Project Architecture and Design}

\subsection{Why the Project Uses a Pipeline Architecture}

The project is implemented as a sequence of connected stages instead of a
single notebook or one large script. This is a deliberate architectural choice.
Because the workflow starts from large raw files and ends with tables, figures,
metrics, and reports, the project benefits from having clear stage boundaries.
Each stage has a defined input and output, and those outputs are persisted
rather than kept only in memory.

This architecture improves three things at once. First, it improves
reproducibility, because downstream stages can be rerun against known upstream
tables. Second, it improves engineering clarity, because the role of each stage
is explicit. Third, it improves presentation quality, because the project can
be explained naturally as a data flow from source files to research outputs.

\subsection{End-to-End Flow}

The full workflow is:

\begin{center}
\fbox{
\begin{minipage}{0.92\textwidth}
\centering
\texttt{Raw BTS .asc files + day.csv} \\
$\Downarrow$ \\
\texttt{clean\_data.py} $\rightarrow$ \texttt{t100\_segment, earnings\_day} \\
$\Downarrow$ \\
\texttt{feature\_engineering.py} $\rightarrow$ \texttt{engineered\_features} \\
$\Downarrow$ \\
\texttt{train\_model.py} $\rightarrow$ model bundle \\
$\Downarrow$ \\
\texttt{backtest\_model.py} $\rightarrow$ \texttt{sde\_predictions, sde\_run\_meta} \\
$\Downarrow$ \\
\texttt{factor\_engineering.py} $\rightarrow$ \texttt{quarterly\_factors} \\
$\Downarrow$ \\
\texttt{fundamental\_backtest.py} $\rightarrow$ portfolio metrics \\
$\Downarrow$ \\
\texttt{plotting.py + agents} $\rightarrow$ charts, run report, anomaly report
\end{minipage}
}
\end{center}

This flow is simple enough to be explained in a classroom setting but strong
enough to support a non-trivial amount of data and analysis. The project is not
trying to imitate a distributed production platform. Instead, it uses a local,
inspectable architecture that is appropriate for a student research project.

\subsection{SQLite as the Data Bus}

The core design choice is the use of SQLite as a central data bus. The database
\texttt{data/quant\_flights.db} sits in the middle of the entire project. Raw
inputs are cleaned into it, intermediate feature tables are written to it,
prediction outputs are persisted to it, and downstream factor and price tables
are also stored there.

This decision has several advantages:

\begin{enumerate}[leftmargin=*, label=(\arabic*)]
  \item SQLite is lightweight and local, which keeps the project portable.
  \item Tables can be inspected directly outside the Python runtime.
  \item Intermediate artifacts survive across runs.
  \item Later stages can reuse existing outputs rather than recompute them.
  \item The whole pipeline becomes easier to debug because each stage leaves a
  trace in a persistent form.
\end{enumerate}

For a project whose cleaned data already exceeds six million rows, this is a
practical and meaningful engineering choice. It turns the codebase into more
than a sequence of calculations; it turns it into a reproducible system.

\subsection{Stage Boundaries and Responsibility}

The architecture also works because each major script has a clear role.
\texttt{clean\_data.py} owns ingestion. \texttt{feature\_engineering.py} owns
feature creation. \texttt{train\_model.py} owns model fitting.
\texttt{backtest\_model.py} owns out-of-sample prediction generation and cache
logic. \texttt{factor\_engineering.py} owns quarterly factor creation.
\texttt{fundamental\_backtest.py} owns the final DS evaluation layer.
\texttt{plotting.py} and the two lightweight agents turn the numeric outputs
into human-readable deliverables.

This script-level separation supports maintainability. It also makes the final
project easier to explain because the architecture is reflected directly in the
file structure.

\subsection{Cached Execution}

One of the most practical features of the architecture is that it supports
cached execution. If a table such as \texttt{engineered\_features} already
exists, or if the prediction signature stored in \texttt{sde\_run\_meta}
matches the current state, the pipeline can skip repeated work. This is
especially valuable for the prediction layer, which would otherwise be
expensive to regenerate every time.

The significance of this feature is not only performance. It also shows that
the project has been engineered as a workflow rather than as a one-time script.
That is an important architectural quality for a coursework submission.

\subsection{Architecture as a DE Contribution}

The architecture itself is one of the project’s main contributions. It is easy
to treat data engineering as merely a preliminary step before “real modeling,”
but in this project the engineering structure is essential to everything that
follows. Without a clean ingestion layer, a persistent storage design, and
cache-aware stage boundaries, the model and backtest layers would be much
harder to trust, rerun, or communicate.

\section{Data Sources and Dataset Links}

\subsection{Main Data Inputs}

The current project uses three main sources of information:

\begin{enumerate}[leftmargin=*, label=(\arabic*)]
  \item BTS T-100 segment data for route-level airline operations;
  \item a local earnings calendar CSV for quarterly event timing; and
  \item cached daily adjusted-close equity prices for the six target airlines
  and SPY benchmark.
\end{enumerate}

These sources are chosen because they match the project objective directly. The
first provides the operating layer, the second provides event alignment, and
the third provides the market layer required for evaluation.

\begin{table}[H]
\centering\small
\caption{Source datasets used by the project}
\label{tab:source_summary}
\begin{tabularx}{\textwidth}{l X X}
\toprule
\textbf{Source} & \textbf{Role in the project} & \textbf{Access / location} \\
\midrule
BTS T-100 Segment data & Route-level monthly operating records including seats, passengers, distance, origin, destination, and carrier & \href{https://www.transtats.bts.gov/}{https://www.transtats.bts.gov/}, stored locally as \texttt{data/*.asc} \\
Local earnings calendar & Quarter-level event dates for the six target airlines & \texttt{data/day.csv}, imported to \texttt{earnings\_day} \\
Adjusted-close price data & Daily equity prices used in the final event-driven backtest & cached in \texttt{equity\_adj\_close\_daily} \\
\bottomrule
\end{tabularx}
\end{table}

\subsection{Why BTS T-100 is the Core Dataset}

The most important dataset in the project is the BTS T-100 segment data. It is
the reason the project is structurally different from a standard price-only
factor exercise. Instead of beginning from market data or third-party analyst
signals, the project starts from the operating layer of the airline business.
This gives the pipeline a strong business logic: passenger flow, seat capacity,
and traffic volume are directly tied to the underlying economics of airline
performance.

Using BTS T-100 also creates a genuine engineering task. The data arrives as
large raw ASCII files rather than in a ready-made model table. That means the
project must solve ingestion, parsing, cleaning, and persistence before
modeling can even begin.

\subsection{Dataset Coverage in the Local Project}

In its current checked form, the project database contains:

\begin{itemize}[leftmargin=*]
  \item 6,724,354 cleaned segment rows in \texttt{t100\_segment};
  \item 6,723,959 engineered rows in \texttt{engineered\_features};
  \item 2,656,668 out-of-sample prediction rows in \texttt{sde\_predictions};
  \item 2,550 quarterly factor rows in \texttt{quarterly\_factors};
  \item 216 earnings-calendar rows in \texttt{earnings\_day}; and
  \item 16,107 cached price rows in \texttt{equity\_adj\_close\_daily}.
\end{itemize}

These counts matter because they show that the project is operating on a
substantial local dataset rather than on a toy subset. The storage and feature
design decisions are therefore motivated by real data scale.

\subsection{Target Airline Universe}

The event-driven backtest focuses on six listed U.S. carriers:

\begin{center}
\texttt{AAL, DAL, UAL, LUV, ALK, JBLU}
\end{center}

This choice keeps the project focused and manageable while still representing a
meaningful set of major U.S. airlines. The operating layer can contain many
more carriers in the raw data, but the final market-facing evaluation is built
around a targeted tradable universe.

\subsection{Benchmark and Event Alignment}

The final backtest uses SPY as benchmark in the cached price layer. This allows
the project to express the final result in an event-driven equity framework
rather than only as a raw operating forecast. The earnings calendar provides
the event dates needed to map the quarterly factor layer onto a later realized
market outcome.

\subsection{Why the Data Design Supports the Report}

The data-source design helps the report because it maps cleanly onto the
narrative:

\begin{itemize}[leftmargin=*]
  \item the operating layer explains \emph{where the signal comes from};
  \item the event calendar explains \emph{when the signal is evaluated}; and
  \item the cached price layer explains \emph{how the final market result is
  measured}.
\end{itemize}

This alignment between source design and report narrative is one reason the
project reads as a coherent DE+DS system instead of a collection of unrelated
scripts.

\section{Data Engineering Implementation}

\subsection{From Raw Files to Structured Tables}

The engineering work begins with raw BTS T-100 files stored locally in
\texttt{data/}. These files are large, historical, and operational in nature.
They are not ready-made model matrices. Instead, they represent a raw source
that must be converted into a structured research asset. This is why the first
substantive contribution of the project is the creation of a stable ingestion
and cleaning pipeline.

\texttt{clean\_data.py} performs this transition. It reads the raw files,
parses the core segment attributes, and writes the result into
\texttt{t100\_segment}. It also imports the local earnings calendar into
\texttt{earnings\_day}. The result is that the project moves immediately from
loosely structured source files to explicit relational tables with predictable
names and roles.

\subsection{Why Table Materialization Matters}

Persisting intermediate tables is one of the strongest engineering decisions in
the project. If the pipeline only transformed raw files into in-memory data
frames, every later stage would require repeated heavy preprocessing and would
be harder to debug. By materializing \texttt{t100\_segment},
\texttt{engineered\_features}, \texttt{sde\_predictions}, and
\texttt{quarterly\_factors}, the project creates a stable intermediate state at
each major boundary.

This has practical benefits:

\begin{itemize}[leftmargin=*]
  \item repeated runs can skip expensive stages;
  \item errors can be isolated to a smaller part of the workflow;
  \item outputs can be inspected directly in the database;
  \item team members can reason about the project stage by stage.
\end{itemize}

Materialization therefore serves both engineering and collaboration purposes.

\subsection{Feature Engineering as a Data Engineering Problem}

Feature engineering is often described purely as a modeling step, but in this
project it is also an engineering task. The feature stage must transform more
than six million segment rows into a model-ready table without exhausting
memory or losing track of provenance. That means the logic must be both
statistically useful and operationally manageable.

The feature pipeline includes:

\begin{itemize}[leftmargin=*]
  \item transformed capacity variables;
  \item route-share measures;
  \item route concentration measures based on HHI;
  \item lagged operating variables;
  \item airport-encoding style features; and
  \item explicit carrier and month structure.
\end{itemize}

These are not arbitrary additions. Each family connects a raw operational idea
to a model-ready quantity while still preserving the semantics of the original
airline data.

\subsection{Memory and Runtime Considerations}

The engineering implementation had to be adjusted so that the pipeline could
run on a local machine with realistic resource limits. This is one reason the
project is more than a conceptual design: it has already been shaped by real
runtime constraints.

The current implementation addresses local execution by:

\begin{enumerate}[leftmargin=*, label=(\arabic*)]
  \item using lighter dtypes in large data frames;
  \item converting suitable columns to categorical form;
  \item avoiding unnecessarily large concatenation patterns in feature and
  prediction matrix construction;
  \item chunking parts of the prediction layer where appropriate; and
  \item relying on cached SQLite outputs instead of repeating every stage.
\end{enumerate}

These changes matter because they transform the pipeline from something that is
theoretically well designed into something that can actually be executed.

\subsection{Feature Table Scale}

The final engineered feature table contains 6,723,959 rows. This is important
because it shows that feature creation was performed at the same order of scale
as the cleaned operational table. The project does not reduce the engineering
problem by discarding most of the data before modeling. Instead, it carries the
operating detail forward into a model-ready layer.

\subsection{Engineering Reuse Through Caching}

Caching appears in the project in two main forms. The first is table existence:
if \texttt{engineered\_features} already exists and the user does not force a
rebuild, the feature stage can be skipped. The second is signature-based
caching in the prediction stage, where \texttt{sde\_run\_meta} determines
whether the prediction layer can be safely reused.

This is a meaningful DE contribution because it links computation to state. The
pipeline is aware of what has already been produced, and therefore behaves like
a workflow system rather than a one-shot script.

\subsection{Engineering Outputs Beyond the Database}

The project’s engineering layer also extends beyond the database itself. Logs
are written to \texttt{logs/pipeline.log}, reports are written to
\texttt{reports/}, and plots are written to \texttt{reports/figures/}. This
means that the engineering design includes both persistent data storage and
persistent communication artifacts.

That is important for a project report because it shows that the engineering
layer does not stop at computation. It also supports observability,
traceability, and presentation.

\subsection{Why DE is a Core Contribution}

Taken together, the ingestion logic, SQLite-centered storage, materialized stage
outputs, caching rules, runtime management, and output organization make data
engineering one of the main contributions of the project. Without this layer,
the later DS results would be much less convincing. The predictive model and
backtest are valuable, but they are valuable \emph{because} they sit on top of
a stable and inspectable engineering foundation.

\section{Data Science Methodology}

\subsection{Modeling Objective}

The DS layer begins once the engineered feature table exists. The main modeling
objective is to estimate airline load factor from route-level operating
features. Load factor is an attractive target because it connects capacity and
demand, and because it can be used as a building block for later passenger and
traffic-related quantities.

This objective is important because it keeps the model grounded in business
structure. Rather than predicting price directly, the model first predicts an
operational quantity that has a natural economic interpretation.

\subsection{XGBoost as the Main Predictive Model}

The project uses an XGBoost regressor as the main predictive engine. This is a
reasonable choice for several reasons. The relationship between route-level
features and load factor is likely nonlinear. Carrier effects, route structure,
seasonality, and competition may all interact. Tree-based boosting is therefore
a practical modeling choice because it can capture such interactions without
requiring a manually specified high-order parametric form.

The final model is saved as \texttt{flight\_revenue\_bundle.joblib}, which is
important from a DS workflow perspective because it preserves the trained
artifact for reuse.

\subsection{Training Window and Validation Discipline}

The current training setup uses historical data up to 2016 for model fitting,
while the out-of-sample evaluation begins in 2017. This is a sensible
time-ordered design for a forecasting-oriented problem. The implementation also
uses a cleaner separation between training and validation during model fitting,
which reduces the risk of optimistic early-stopping behavior.

In addition, the residual or uncertainty calibration used downstream is
constructed in a walk-forward manner. This detail matters because the project
is not merely reporting in-sample fit quality; it is trying to produce an
operationally meaningful out-of-sample layer.

\subsection{Prediction Layer and Derived Quantities}

Once the model is trained, it is applied to the out-of-sample period from 2017
to 2025. The prediction layer generates quantities such as predicted passengers
and predicted $\RPK$, and stores them in \texttt{sde\_predictions}. The table
also retains the realized passenger and traffic information required for later
diagnostics.

In project terminology this stage is described as an SDE backtest layer, but in
practical terms it functions as a stochastic operating-forecast layer that
extends the model output beyond a pure point estimate. The purpose is to create
a richer prediction object that can later be aggregated into quarterly
research factors.

\subsection{Prediction Metrics}

The current checked project snapshot reports:

\begin{itemize}[leftmargin=*]
  \item MAE of approximately 281.30;
  \item RMSE of approximately 640.90;
  \item $\WMAPE$ of approximately $11.34\%$; and
  \item 2,656,668 out-of-sample prediction observations.
\end{itemize}

These metrics indicate that the route-level forecast layer is stable enough to
support the later factor and backtest stages. They do not need to be perfect to
be useful. What matters here is that the prediction layer provides a coherent
operational signal that can be transformed into carrier-level quarterly factors.

\begin{figure}[H]
\centering
\includegraphics[width=0.88\textwidth]{sde_pred_vs_actual_scatter_20260405_225923.png}
\caption{Predicted versus actual passengers in the out-of-sample prediction
layer. The figure summarizes how closely the model-based forecast aligns with
the realized operating outcome.}
\label{fig:sde_scatter}
\end{figure}

\begin{figure}[H]
\centering
\includegraphics[width=0.88\textwidth]{carrier_wmape_bar_20260405_225924.png}
\caption{Carrier-level WMAPE summary. This chart is useful because it turns the
route-level forecast problem into a carrier-level diagnostic view.}
\label{fig:wmape_bar}
\end{figure}

\subsection{Quarterly Factor Construction}

The monthly prediction outputs are not the final research object. They are
aggregated into quarterly carrier-level factors that summarize operating
conditions in a form suitable for cross-sectional event analysis. The project
currently uses four factors:

\begin{center}
\texttt{factor\_rpk, factor\_lf, factor\_ask, factor\_share}
\end{center}

These factors capture different aspects of airline operations. RPK reflects
traffic, LF reflects utilization, ASK reflects capacity, and share-related
measures reflect route position or competitive structure. Together they form a
small but interpretable factor family.

\subsection{Why the DS Layer is Sequential}

One of the strengths of the DS implementation is that it is sequential and
interpretable:

\begin{enumerate}[leftmargin=*, label=(\arabic*)]
  \item route-level features are built from operational records;
  \item those features feed a predictive model;
  \item model outputs are stored at the route and month level;
  \item stored outputs are aggregated into quarterly carrier factors; and
  \item the factors are evaluated in an event-driven backtest.
\end{enumerate}

This sequence gives the DS layer more depth than a standalone prediction score.
The project does not stop when the model finishes; it carries the output
forward until it becomes a market-facing research result.

\section{Backtest Design, Results, and Interpretation}

\subsection{Why the Final Evaluation is Event Driven}

The final research question of the project is not whether airline operations
can be forecast in isolation. It is whether those operational signals can be
organized into an equity-event framework around earnings releases. This is why
the last stage of the DS layer is event driven.

The final backtest uses the quarterly factor table, the earnings calendar, and
cached equity prices for the six target airlines and SPY benchmark. It asks
whether carriers with stronger operating signals subsequently show stronger
event-window outcomes.

\subsection{Tradable Post-Formation Return Design}

An important part of the current implementation is that the event return is
measured in a way that matches the portfolio formation logic. The signal is
formed after quarter end, and the measured return occurs in a post-formation
event window. This makes the evaluation more practically aligned and easier to
defend in a report setting.

This design matters because the credibility of a backtest depends not only on
the final metric but also on whether the timing logic is internally consistent.

\subsection{MWU Expert Combination}

The project combines the four quarterly factors using a multiplicative-weights
update framework. In simple terms, each factor acts as an expert that ranks the
cross-section, and the system adapts the expert weights through time based on
realized outcomes. This keeps the final backtest interpretable while still
allowing the combination rule to respond to observed performance.

For a coursework project, this is a good compromise between simplicity and
adaptivity. The combination logic is more interesting than a fixed equal-weight
average, but it is still understandable enough to explain clearly in a report
or presentation.

\subsection{Headline Results}

The current checked snapshot from \texttt{reports/\_last\_run\_metrics.json}
reports:

\begin{table}[H]
\centering\small
\caption{Headline checked results}
\label{tab:headline_results}
\begin{tabular}{lr}
\toprule
\textbf{Metric} & \textbf{Value} \\
\midrule
Prediction $\WMAPE$ & 11.34\% \\
Prediction observations & 2,656,668 \\
Backtest quarters & 31 \\
Gross cumulative return & 91.47\% \\
Net cumulative return & 68.15\% \\
Gross annualized Sharpe & 1.287 \\
Net annualized Sharpe & 1.027 \\
Quarterly win rate & 67.74\% \\
\bottomrule
\end{tabular}
\end{table}

These metrics support a positive reading of the final system. The operating
prediction layer is stable at scale, and the event-driven portfolio layer
remains positive even after transaction-cost adjustment. That does not imply
that every factor is individually strong in every quarter, but it does show
that the pipeline produces a coherent DS output rather than only an isolated
engineering artifact.

\subsection{Portfolio-Level Visualization}

\begin{figure}[H]
\centering
\includegraphics[width=0.88\textwidth]{quarterly_cumulative_return_gross_net_20260405_225924.png}
\caption{Quarterly cumulative gross and net portfolio returns. The net series
remains positive, which is important for interpreting the strategy beyond a raw
gross-return view.}
\label{fig:gross_net}
\end{figure}

The cumulative return chart is one of the clearest visual summaries in the
project because it compresses the final event-driven evaluation into a single
readable figure. It shows that the gross and net paths differ, but both remain
positive in the checked run.

\subsection{Factor Interpretation}

The factor layer is summarized in the mean information coefficients. In the
current run, the ASK and RPK related factors are stronger than the LF factor.
This is a sensible outcome because demand scale and capacity structure often
carry more cross-sectional business information than utilization alone.

\begin{figure}[H]
\centering
\includegraphics[width=0.78\textwidth]{factor_ic_bar_20260405_225924.png}
\caption{Average factor information coefficients from the current checked run.}
\label{fig:factor_ic2}
\end{figure}

\begin{figure}[H]
\centering
\includegraphics[width=0.78\textwidth]{factor_corr_heatmap_20260405_225924.png}
\caption{Factor correlation heatmap. This chart helps show where factor overlap
exists and where the signal family may still provide differentiated views.}
\label{fig:factor_corr2}
\end{figure}

The factor interpretation should be balanced. The project does not claim that
each factor is independently sufficient. Instead, it shows that a small,
operationally grounded factor family can be organized into an interpretable
backtest framework.

\subsection{Regime Interpretation}

\begin{figure}[H]
\centering
\includegraphics[width=0.76\textwidth]{regime_performance_bar_20260405_225924.png}
\caption{Regime performance summary from the checked snapshot.}
\label{fig:regime2}
\end{figure}

The regime breakdown is useful because it shows that the project is not being
judged only on one aggregate number. The checked results remain positive across
pre-COVID, COVID, and post-COVID groupings, with the strongest Sharpe appearing
in the post-COVID segment. This helps the report present the system as a
structured research framework with interpretable behavior across broad periods.

\subsection{What the Results Mean}

The most reasonable interpretation of the current DS results is that the
pipeline is directionally successful. The route-level model produces stable
operational forecasts, those forecasts can be converted into quarterly factors,
and the resulting factor framework produces positive event-driven portfolio
results in the checked run.

This should not be oversold as a final trading product. Rather, it should be
presented as a coherent research outcome: a DE+DS system that demonstrates the
practical value of organizing airline operating data into a market-evaluation
framework.

\section{Execution Instructions, Outputs, and Reproducibility}

\subsection{Environment and Dependencies}

The project is designed to run locally in a Python environment with the
packages listed in \texttt{requirements.txt}. The core dependencies include
pandas, numpy, xgboost, scipy, matplotlib, yfinance, PyYAML, joblib, and the
OpenAI Python client used by the lightweight reporting tools.

The local execution design is important because it aligns with the architecture
described earlier. The project does not require a remote warehouse, cluster, or
orchestration service. A user with the required files and packages can run the
pipeline from the project directory and inspect all outputs locally.

\subsection{Core Commands}

The most relevant commands for reproducing the workflow are:

\begin{lstlisting}[language=bash, caption={Practical command set for the project}]
python -m pip install -r requirements.txt
python clean_data.py
python main.py

# rebuild variants
python main.py --force-features --force-retrain --force-sde --force-factors
python main.py --force-retrain --force-sde --force-factors
python main.py --force-sde --force-factors
\end{lstlisting}

These commands are intentionally simple. For a coursework report, that matters:
the execution path is easy to document and easy to demonstrate.

\subsection{Role of \texttt{config.yaml}}

The configuration file is one of the project’s strongest software-engineering
features. It groups settings into the same structure as the pipeline itself:
data paths, model settings, backtest settings, factor settings, fundamental
evaluation settings, plotting settings, logging, and agent controls.

This design helps both reproducibility and explainability. A reader can inspect
\texttt{config.yaml} and understand the operational assumptions of the run
without searching through multiple Python files.

\subsection{Expected Outputs}

After a successful run, the user should expect the following artifacts:

\begin{itemize}[leftmargin=*]
  \item updated tables in \texttt{data/quant\_flights.db};
  \item a model bundle in \texttt{models/};
  \item a run report in \texttt{reports/};
  \item a metrics snapshot in \texttt{reports/\_last\_run\_metrics.json};
  \item an anomaly report in \texttt{logs/}; and
  \item six PNG charts in \texttt{reports/figures/}.
\end{itemize}

This output set is one of the reasons the project reads as an integrated
workflow. It produces persistent engineering outputs, analytical outputs, and
presentation outputs from a single orchestrated execution.

\subsection{Current-Run Automation}

The current project contains two lightweight supporting tools:

\begin{enumerate}[leftmargin=*, label=(\roman*)]
  \item a run report agent that summarizes the current project outputs; and
  \item an anomaly agent that summarizes current-run quality signals.
\end{enumerate}

These agents are intentionally limited. They do not drive the pipeline and they
are not positioned as the main contribution of the project. Their role is to
improve usability by translating the current state of the pipeline into a short
human-readable layer.

\subsection{Reproducibility Perspective}

From a reproducibility standpoint, the project is well structured for a local
coursework pipeline:

\begin{itemize}[leftmargin=*]
  \item source files are kept locally;
  \item main intermediate outputs are persisted to SQLite;
  \item model artifacts are saved by name;
  \item run metrics are exported in machine-readable form;
  \item figures are written to a predictable output folder;
  \item the main execution log records stage progression.
\end{itemize}

This combination makes the project easier to rerun and easier to defend in a
report setting, because the system state is not hidden inside a transient
runtime.

\section{Findings, Future Work, and Conclusion}

\subsection{Main Findings}

The main finding of the project is that airline operating data can be organized
into a coherent DE+DS workflow that produces interpretable and positive
end-to-end research outputs. The engineering layer succeeds in turning raw BTS
files into reusable tables and artifacts. The modeling layer succeeds in
turning those tables into a forecast-and-factor system that can be evaluated in
an event-driven equity framework.

At the prediction level, the checked $\WMAPE$ of approximately 11.34\% suggests
that the operational model is stable enough to support later analysis. At the
portfolio level, the checked gross and net cumulative returns and Sharpe ratios
show that the signal framework remains positive even after cost adjustment.

\subsection{Why the DE Story Matters}

It would be easy to read the project only through the final backtest chart, but
that would understate its strongest technical contribution. The reason the
system works as a project is that the engineering layer is explicit, persistent,
and inspectable. The workflow begins with raw files, passes through named
tables, stores model and prediction outputs, and ends with charts and reports.

That design gives the later DS results more credibility because they are
produced by a stable workflow rather than by an opaque one-off script.

\subsection{Why the DS Story Matters}

The DS contribution is equally important. The project does not stop at route
forecasting. It carries the signal forward into quarterly factor construction
and event-driven market evaluation. That sequence gives the final result a much
stronger research interpretation than a route-level model score alone.

The checked figures and metrics show that the DS layer is not only technically
implemented but also communicable. The model produces a readable scatter plot,
carrier-level WMAPE diagnostics, factor comparison charts, regime charts, and a
final gross-versus-net portfolio view.

\subsection{Future Work}

Several future directions would strengthen the project further:

\begin{enumerate}[leftmargin=*, label=\textbf{F\arabic*.}]
  \item \textbf{Broader airline universe.} Expanding the cross-section would
  deepen the factor layer and support richer evaluation.

  \item \textbf{Additional external features.} Fuel proxies, macro conditions,
  and richer network variables could be added to the current operational set.

  \item \textbf{Extended validation.} Additional event-window and transaction
  cost sensitivity analysis would improve the robustness of interpretation.

  \item \textbf{More formal test coverage.} Additional unit and integration
  tests would strengthen maintainability and handoff quality.
\end{enumerate}

These future directions do not undermine the current system. Instead, they
represent natural next steps for extending a pipeline that already has a solid
engineering and analytical core.

\subsection{Final Conclusion}

QF5214 should be understood as a data-engineering-driven quantitative research
project. Its core achievement is the combination of structured ingestion,
feature construction, predictive modeling, factor generation, event-driven
evaluation, and automated output generation inside one coherent local workflow.

For a student group project, this is a meaningful outcome. The system shows
that data engineering and data science can be integrated into a practical
research pipeline with clear storage design, interpretable analysis, and
presentable outputs. It is this end-to-end quality, rather than any isolated
single model score, that defines the project’s value.

\appendix

\section{Appendix: Technical Reference}

\subsection{Core Database Tables}

\begin{longtable}{p{3.2cm} p{2.0cm} p{8.9cm}}
\caption{Core database table reference}\label{tab:appendix_tables}\\
\toprule
\textbf{Table} & \textbf{Rows} & \textbf{Meaning} \\
\midrule
\endfirsthead
\toprule
\textbf{Table} & \textbf{Rows} & \textbf{Meaning} \\
\midrule
\endhead
\texttt{t100\_segment} & 6,724,354 & Cleaned route-level operational base table derived from raw BTS files. \\
\texttt{earnings\_day} & 216 & Quarterly event-date reference for the six airline tickers. \\
\texttt{engineered\_features} & 6,723,959 & Route-level modeling table containing transformed operational features. \\
\texttt{sde\_predictions} & 2,656,668 & Out-of-sample route-level prediction and derived traffic output table. \\
\texttt{sde\_run\_meta} & 1 & Metadata used to determine whether prediction outputs can be reused. \\
\texttt{quarterly\_factors} & 2,550 & Airline-quarter factor table consumed by the final backtest. \\
\texttt{equity\_adj\_close\_daily} & 16,107 & Cached adjusted-close equity prices for six airlines and SPY. \\
\bottomrule
\end{longtable}

\subsection{Main Script Responsibilities}

\begin{longtable}{p{3.2cm} p{9.6cm}}
\caption{Script-level responsibility map}\label{tab:script_map}\\
\toprule
\textbf{Script} & \textbf{Role in the workflow} \\
\midrule
\endfirsthead
\toprule
\textbf{Script} & \textbf{Role in the workflow} \\
\midrule
\endhead
\texttt{clean\_data.py} & Parses and cleans raw BTS operating files and imports the local earnings calendar. \\
\texttt{feature\_engineering.py} & Builds route-level model features from the cleaned operational base table. \\
\texttt{train\_model.py} & Trains the XGBoost load-factor model and stores the model bundle. \\
\texttt{backtest\_model.py} & Generates the out-of-sample prediction layer and manages cache metadata. \\
\texttt{factor\_engineering.py} & Aggregates monthly prediction outputs into quarterly airline factors. \\
\texttt{fundamental\_backtest.py} & Runs the final event-driven MWU backtest. \\
\texttt{plotting.py} & Exports current-run charts for interpretation and presentation. \\
\texttt{main.py} & Orchestrates the full pipeline. \\
\bottomrule
\end{longtable}

\subsection{Current Checked Metrics Snapshot}

\begin{longtable}{p{6.0cm} p{4.8cm}}
\caption{Current checked metrics snapshot}\label{tab:metrics_snapshot}\\
\toprule
\textbf{Metric} & \textbf{Value} \\
\midrule
\endfirsthead
\toprule
\textbf{Metric} & \textbf{Value} \\
\midrule
\endhead
Prediction MAE & 281.2983 \\
Prediction RMSE & 640.9028 \\
Prediction WMAPE & 0.1134 \\
Prediction sample size & 2,656,668 \\
Gross cumulative return & 0.9147 \\
Net cumulative return & 0.6815 \\
Gross annualized Sharpe & 1.2871 \\
Net annualized Sharpe & 1.0273 \\
Quarterly win rate & 0.6774 \\
Mean IC: factor\_ask & 0.0912 \\
Mean IC: factor\_rpk & 0.0802 \\
Mean IC: factor\_share & 0.0740 \\
Mean IC: factor\_lf & 0.0359 \\
\bottomrule
\end{longtable}

\subsection{Figure Set Used in the Report}

The current checked figure set used by the report consists of:

\begin{enumerate}[leftmargin=*, label=(\arabic*)]
  \item predicted versus actual passengers;
  \item carrier-level WMAPE bar chart;
  \item cumulative gross and net return chart;
  \item factor IC bar chart;
  \item factor correlation heatmap; and
  \item regime performance chart.
\end{enumerate}

These figures span the full analytical pipeline: model fit, carrier diagnostic,
portfolio path, factor ranking, factor interaction, and regime interpretation.

\subsection{Suggested Demonstration Sequence}

If the project is presented live, a concise demonstration sequence would be:

\begin{enumerate}[leftmargin=*, label=(\arabic*)]
  \item show the project folders and local data assets;
  \item explain the SQLite-centered stage architecture;
  \item run or describe \texttt{python main.py};
  \item inspect the latest metrics snapshot;
  \item open the key figures;
  \item conclude with the DE and DS contributions.
\end{enumerate}

This sequence follows the same logic as the report itself, making the written
document and oral presentation mutually reinforcing.

\subsection{Representative Feature Logic}

The engineered feature layer is central enough to deserve a more explicit
technical summary. Although the report does not list every derived column, the
main logic can be grouped into the following representative transformations:

\begin{enumerate}[leftmargin=*, label=(\arabic*)]
  \item \textbf{volume transformation:} route-level seat and passenger values
  are transformed so that scale can be represented smoothly rather than only
  through raw counts;
  \item \textbf{market structure:} route share and HHI-like quantities are
  created to describe whether a route is highly concentrated or relatively more
  competitive;
  \item \textbf{seasonal memory:} lagged values are introduced so that the
  model can compare current route activity with earlier route behavior,
  especially the same period in a prior annual cycle;
  \item \textbf{location encoding:} origin and destination airports are mapped
  into compact representations that preserve airport effects without exploding
  the dimensionality of the feature space;
  \item \textbf{carrier and calendar structure:} carrier identity and month
  information provide an explicit structure for systematic differences across
  airlines and within-year traffic patterns.
\end{enumerate}

These feature groups are not meant to impress through quantity. They matter
because they encode airline economics in a way that remains interpretable. The
feature layer therefore supports the project’s DS goal while remaining faithful
to its DE foundation.

\subsection{Why the Feature Layer is Not Arbitrary}

A common weakness in student modeling projects is that features are added
without a clear story. This project avoids that problem by keeping feature
families tied to recognizable business concepts. Capacity variables correspond
to supply, passenger and RPK-related variables correspond to realized demand,
share and HHI correspond to competition, and lag features correspond to
historical continuity and seasonality.

This matters because the report can therefore explain the feature space in
domain language rather than only in code language. That improves both the
analytical quality of the project and the persuasiveness of the final written
report.

\subsection{Expanded Backtest Reading Notes}

The final backtest should be read with a few practical interpretations in mind.
First, the portfolio output is the end result of the full DE+DS chain rather
than the output of one isolated formula. Second, the gross and net return paths
should be interpreted together. Gross performance shows the directional signal,
while net performance shows whether that signal survives a simple implementation
friction assumption. Third, the factor IC charts should be interpreted as a
compact way to compare factor usefulness, not as a claim that every factor is
independently sufficient.

These distinctions are helpful because they encourage a more serious reading of
the project. The report is not claiming a magic factor. It is showing that a
structured operating-data pipeline can support a layered research process from
raw data to market evaluation.

\subsection{Interpreting the Current Performance Mix}

The checked metrics suggest a project whose pieces are aligned. The route-level
prediction layer is strong enough to generate meaningful downstream inputs. The
quarterly factor layer is interpretable and non-trivial. The final backtest is
positive both gross and net. The charts reinforce that the results can be
understood visually rather than only through one table of numbers.

This mixture is important. In many small projects, either the engineering is
strong but the analytics remain shallow, or the analytics are ambitious but the
engineering remains fragile. The current QF5214 version is notable because it
connects the two layers in a visible way.

\subsection{Operational Value of the Logs and Reports}

The project’s logs and reports also deserve explicit mention as engineering
artifacts. The pipeline log records which stages were skipped, which were run,
how long they took, and whether cache logic was used. The markdown run report
records the current checked metrics in narrative form. The anomaly report
records current-run issues in a concise diagnostic format.

These outputs are not the core analytical result, but they improve the quality
of the workflow. They turn a technical pipeline into a project that can be
reviewed by other people more easily. In a team setting, this matters because
it reduces the amount of hidden context required to understand the latest run.

\subsection{Why Lightweight Agents Still Make Sense}

The report intentionally does not place the agent layer at the center of the
project. Even so, it is still worth explaining why these tools are useful. The
two current agents are practical because they work on top of the pipeline
rather than trying to replace it. The report agent summarizes the latest
outputs. The anomaly agent points out obvious current-run concerns. Neither
agent changes the computational core or claims to discover the signal by
itself.

This restrained use of automation fits the overall philosophy of the project.
The engineering and modeling layers do the heavy lifting; the automation layer
improves readability and presentation.

\subsection{Extended Reproducibility Notes}

Reproducibility in this project is not defined only by whether the same Python
script can be run twice. It is also defined by whether each major stage leaves
behind a persistent and inspectable result. The local SQLite database preserves
the operational base table, the feature table, the prediction outputs, the
factor table, and the cached prices. The model bundle preserves the trained
predictive artifact. The figures preserve the visual analytical layer. The log
and markdown outputs preserve the narrative and runtime interpretation layer.

This multi-layer reproducibility is one of the project’s strongest qualities.
It is also why the report is justified in emphasizing data engineering so
heavily: the reproducibility of the whole system depends more on the structure
of the workflow than on any one formula.

\subsection{Suggested Division of Work in a Team Context}

The project structure also maps naturally onto team roles, which is useful for
explaining it as a group assignment. One team member could focus on raw data
handling and SQLite storage. Another could focus on feature engineering and the
predictive model. Another could focus on the quarterly factor layer and
backtest. Another could focus on presentation outputs such as charts and final
documentation. Although these responsibilities overlap in practice, the
architecture supports this division because the workflow already has clear stage
boundaries.

This is another sign that the project is well organized. A system with good
stage separation is easier both to build and to present as collaborative work.

\subsection{How the Report Itself Mirrors the Pipeline}

The report is structured to mirror the pipeline’s logic. It begins with the
problem, moves into architecture and source design, then explains data
engineering, then modeling, then evaluation, then execution and outputs, and
finally closes with findings and future work. This ordering is intentional. A
pipeline project is best explained as a chain of transformations rather than as
an unordered list of tools.

That is why the report gives so much attention to stage transitions. The value
of the system lies not only in the existence of each stage, but in how those
stages connect.

\subsection{Expanded Summary of Project Value}

The overall value of QF5214 can be summarized through four statements:

\begin{enumerate}[leftmargin=*, label=(\alph*)]
  \item it demonstrates that raw airline operating data can be organized into a
  durable research database;
  \item it demonstrates that engineered airline features can support
  out-of-sample operational forecasting at scale;
  \item it demonstrates that those forecasts can be turned into quarterly
  factors and evaluated in an earnings-event setting; and
  \item it demonstrates that the full process can be presented through logs,
  charts, and reports rather than only through raw code.
\end{enumerate}

For a student group project, this is a strong combined outcome. It reflects not
only coding effort, but also workflow design, analytical framing, and technical
communication.

\subsection{Closing Appendix Note}

This appendix has documented the project at a level detailed enough for a
reviewer to understand its technical composition even without executing the
code immediately. That is an appropriate finishing point for the report. The
main body explains the story of the project, while the appendix records the
technical reference details that support that story.

\subsection{Final Technical Takeaway}

Seen as a whole, the project can be read as a three-layer system. The first
layer is data engineering: ingesting, structuring, storing, and reusing
airline operating data. The second layer is data science: modeling load factor,
building quarterly factors, and evaluating an event-driven portfolio. The third
layer is communication: logs, charts, reports, and lightweight anomaly
summaries. The reason the project works is that these three layers are aligned
rather than disconnected.

This alignment is the report’s main message. The project is not simply a set of
Python files, and it is not simply a machine-learning exercise. It is a local,
reproducible research workflow that shows how engineering structure can support
analytical quality. For a coursework submission, that is a meaningful and
defensible outcome.

\subsection{Final Presentation Summary}

If the project must be summarized in a final two-minute statement, the summary
would be:

\begin{quote}
QF5214 builds a complete airline quant research workflow from raw BTS operating
data to event-driven equity results. Its main contribution is the combination
of a SQLite-centered data engineering pipeline with an interpretable data
science layer for forecasting, factor construction, and backtesting. The final
system produces persistent data artifacts, readable metrics, and presentation
quality charts, making it a strong example of an end-to-end DE+DS coursework
project.
\end{quote}

\end{document}
